\documentclass[11pt]{article}
\usepackage[T1]{fontenc}
\usepackage[utf8]{inputenc}
\usepackage{lmodern}
\usepackage[english]{babel}
\usepackage[margin=1in]{geometry}
\usepackage{amsmath,amssymb,amsthm,booktabs,array}
\usepackage{xurl}
\usepackage[colorlinks=true,linkcolor=black,citecolor=blue,urlcolor=blue]{hyperref}
\setlength{\emergencystretch}{2em}
\hypersetup{pdftitle={Nine-Square Arithmetic-Derivative Tuples: An Extension of OEIS A249110 and Their Scaling Structure},pdfauthor={Carlo Corti}}
\newcommand{\D}{\operatorname{D}}
\newcommand{\Lrate}{\operatorname{L}}
\newcommand{\NN}{\mathbb{N}}
\newcommand{\PP}{\mathbb{P}}
\newtheorem{theorem}{Theorem}[section]
\newtheorem{proposition}[theorem]{Proposition}
\newtheorem{lemma}[theorem]{Lemma}
\newtheorem{corollary}[theorem]{Corollary}
\theoremstyle{definition}
\newtheorem{definition}[theorem]{Definition}
\theoremstyle{remark}
\newtheorem{remark}[theorem]{Remark}
\title{Nine-Square Arithmetic-Derivative Tuples:\\An Extension of OEIS A249110 and Their Scaling Structure}
\author{Carlo Corti\\Independent researcher\\\texttt{carlo.corti@outlook.com}}
\date{October 2026}
\begin{document}
\maketitle

\begin{abstract}
OEIS A249110 lists the integers whose squares, together with the squares
of their first eight arithmetic derivatives, sum to an integer square.
Starting from the accepted five-term prefix, a segmented C computation
classifies $823544\le n\le22000000000$ and establishes four consecutive
terms, ending at $a(9)=10361580899$. The exclusion claim depends on the
AI code-reviewed producer, its finite-domain primality premises, the complete
records, and the inherited baseline; it is not supported by a second
independent enumeration of all negatives. All nine positive witnesses
were independently reconstructed. The resulting next-term bound is
$22000000000<a(10)\le285311670611=11^{11}$.
The nine ordered ten-coordinate witnesses have five primitive forms.
Classical logarithmic-derivative theory and a product-rule induction
give the complete positive integer scales of each observed form:
two classes are infinite and three are singletons. Besides the known
stationary family $p^p$, the nonconstant form shared by $a(7)$ and
$a(8)$ yields $330004229p^p$ for every prime $p$.
A previously known terminating witness acquires rank six in this
sequence; its zero tail explains persistence across later depths.
The scale results classify the represented forms, not all possible
forms, and no density model is inferred from nine terms.
\end{abstract}

\noindent\textit{2020 Mathematics Subject Classification:}
Primary 11A25; Secondary 11D09, 11Y55.

\noindent\textit{Key words and phrases:}
arithmetic derivative, Pythagorean tuple, integer sequence, primitive
ordered form, logarithmic derivative, exact enumeration.

\section{Definition, prior knowledge and the extension}\label{sec:intro}

Let $\NN=\{0,1,2,\ldots\}$ and let $\PP$ be the primes. The ordinary
arithmetic derivative on $\NN$ is defined by
\begin{equation}\label{eq:derivative}
 \D(0)=\D(1)=0,\qquad
 \D(n)=\sum_{p^e\parallel n}e\frac{n}{p}\quad(n\ge2).
\end{equation}
Here $p^e\parallel n$ means that $p^e$ is the exact power of $p$ dividing
$n$. Equivalently, $\D(p)=1$ for primes and
\begin{equation}\label{eq:product}
 \D(ab)=a\D(b)+b\D(a)\qquad(a,b\in\NN).
\end{equation}
These definitions and their logarithmic formulation are classical
\cite{barbeau,ufnarovski,kovic}. Superscripts on $\D$ denote composition,
with $\D^0(n)=n$. The map is not additive: $\D(2+3)=1$, whereas
$\D(2)+\D(3)=2$.

\begin{definition}\label{def:sequence}
For integers $n,h\ge0$, put
\[
 Q_h(n)=\sum_{j=0}^{h}(\D^j(n))^2.
\]
For $h\ge1$, let
$S_h=\{n\in\NN:n\ge2,\ Q_h(n)\text{ is a square}\}$.
The increasing enumeration of $S_8$ is denoted $a(1),a(2),\ldots$,
with offset one. For $n\in S_8$, define its ordered witness
\begin{equation}\label{eq:witness}
 W(n)=(x_0,x_1,\ldots,x_8,R),\qquad
 x_j=\D^j(n),\quad R>0,\quad \sum_{j=0}^8x_j^2=R^2.
\end{equation}
\end{definition}
Thus the OEIS title's ten coordinates are the nine chain values and
their root; the sequence index is not a coordinate. No primitivity, increasing-leg
order, family membership or strict positivity of later iterates is
required. Later zero coordinates are admitted. Inputs zero and one are
excluded even though the bare square equation has degenerate solutions
there. This is the explicit zero-tail convention used throughout.

The sequence is OEIS A249110, contributed by Paolo P. Lava on
October 21, 2014~\cite{sequence}. The snapshot acquired on October 24,
2026, revision 24 dated February 18, 2021, has the five indexed terms
\begin{equation}\label{eq:baseline}
 4,\quad27,\quad3125,\quad398747,\quad823543.
\end{equation}
The entry checked again on October 27 still displays this prefix.
Its indexing is inherited here: no new full scan below $823544$ was
performed. The Maple call printed in the entry is a program, not a
retained receipt of an exhaustive search. Its predicate imposes no
positivity test on later legs. Moreover, its derivative expression uses
the factor list from \texttt{ifactors}: the documented lists for zero
and one are empty~\cite{maple}, giving zero derivatives. This is a
static compatibility check for the zero-tail convention, not a claim
to have executed Maple or recovered a historical search receipt.

The subsequence $p^p$, for $p\in\PP$, was credited in A249110 to
Ray Chandler on December 22, 2016~\cite{sequence,fixedpoints}.
Its mechanism is the known identity $\D(p^p)=p^p$; it is not a
discovery of the present computation. Also known before this project was
$N=1072281223$, with a terminating chain and a square sum at every
depth at least five, from the author's A249107 study~\cite{corti7}.
The first A249110 campaign determines the rank of this known witness,
namely $a(6)=N$. It does not discover the integer or its zero-tail
mechanism anew.

The first contribution is the finite extension: the subsequent validated
prefix of a continuation supplies ranks seven through nine, and a final
adjacent run establishes an empty interval ending at $22000000000$.
Sections~\ref{sec:method} and~\ref{sec:results} give its method and
complete statement after the structural analysis. The mathematical analysis then has
two distinct inputs: positive chains, whose identities are individually
checkable, and interval coverage, needed to assign consecutive ranks.
Neither input substitutes for the other.

The second contribution is the scale structure of these nine witnesses.
Their normalization gives five forms, including one nonconstant form
represented at scales four and twenty-seven. This relation was recorded
while preparing the manuscript handoff, after the numerical results.
Here it is extended by proof to an infinite family and placed within
the complete scale classification of the five observed forms.
The general principles follow logarithmic-derivative theory
\cite{ufnarovski,kovic} and the author's earlier fixed-form work
\cite{corti2,corti5,corti6,corti7}; these earlier depth-two through
depth-five applications already include nonconstant families. The
depth-eight result concerns the particular form displayed here; no
absolute priority is claimed.

\section{Scaling a chain without an additive derivative}\label{sec:scales}

For positive integers $s$, let
\begin{equation}\label{eq:rate}
 \Lrate(s)=\frac{\D(s)}{s}
   =\sum_{p^e\parallel s}\frac{e}{p},\qquad \Lrate(1)=0.
\end{equation}
Equation~\eqref{eq:product} gives
$\Lrate(st)=\Lrate(s)+\Lrate(t)$ without a coprimality hypothesis.
This is additivity on products, not on sums.
For rational $\rho\ge0$, define
\begin{equation}\label{eq:shift}
 T_\rho(m)=\D(m)+\rho m
\end{equation}
whenever $m$ is a nonnegative integer. An iterate is used only when
every preceding value is a nonnegative integer; no derivative of a
noninteger is needed in this paper.

\begin{lemma}[Product-rule induction]\label{lem:scale}
If $t\ge1$ satisfies $\D(t)=\rho t$ and
$m,T_\rho(m),\ldots,T_\rho^j(m)$ are nonnegative integers, then
\begin{equation}\label{eq:iterate}
 \D^j(tm)=tT_\rho^j(m).
\end{equation}
Consequently, when the iterates through eight are integral,
\begin{equation}\label{eq:Qscale}
 Q_8(tm)=t^2\sum_{j=0}^8(T_\rho^j(m))^2.
\end{equation}
\end{lemma}
\begin{proof}
The claim at $j=0$ is the identity $tm=tm$. If $z=T_\rho^j(m)$ is the
current integer value, the product rule gives
$\D(tz)=t\D(z)+z\D(t)=t(\D(z)+\rho z)=tT_\rho(z)$.
This proves the next step, including $z=0$. Squaring and summing gives
\eqref{eq:Qscale}. No coprimality assumption is used, and no binomial
expansion of an additive operator is invoked.
\end{proof}

\begin{definition}\label{def:form}
For a witness $W(n)$, let $g$ be the gcd of all ten coordinates,
including $R$. Its primitive ordered form is $F=W(n)/g$.
A nonnegative integer tuple
$F=(u_0,\ldots,u_8,r)$ with gcd one, $u_0,r>0$, is called admissible
if $sF=W(su_0)$ for at least one positive integer $s$ with $su_0\ge2$.
Its homothetic class consists of all such witnesses $sF$.
\end{definition}
Normalization preserves order and never includes the index $i$.
An admissible form need not itself be an ordinary derivative chain.

\begin{theorem}[Criterion for a fixed form]\label{thm:criterion}
Let $F=(u_0,\ldots,u_8,r)$ be admissible. Set
\begin{equation}\label{eq:lambda}
 \lambda=\frac{u_1-\D(u_0)}{u_0}.
\end{equation}
Then $\lambda\ge0$,
\begin{equation}\label{eq:transitions}
 u_{j+1}=\D(u_j)+\lambda u_j\ (0\le j<8),
 \qquad \sum_{j=0}^8u_j^2=r^2.
\end{equation}
For every positive integer $s$ with $su_0\ge2$,
$sF=W(su_0)$ if and only if $\Lrate(s)=\lambda$.
Every positive integer scale satisfying $\Lrate(s)=\lambda$
automatically has $su_0\ge2$. No positive noninteger rational multiple of $F$
has all ten coordinates integral.
\end{theorem}
\begin{proof}
Choose a witnessing scale $g$. Dividing
$\D(gu_0)=gu_1$ by $g$ gives
$u_1=\D(u_0)+u_0\Lrate(g)$, so $\lambda=\Lrate(g)\ge0$.
The same argument at every later coordinate gives the transitions,
including zeros. Division of the square identity by $g^2$ gives the
last equation. Necessity for any other scale follows from the first
transition. Conversely, the product rule and
\eqref{eq:transitions} reconstruct all eight derivatives at a scale
with $\Lrate(s)=\lambda$, and the square identity scales by $s^2$.
The input cannot be one: if $s=u_0=1$ and $\Lrate(s)=\lambda$, then
$\lambda=0$, the transitions force $u_1=\cdots=u_8=0$, and the
same rate gives $\Lrate(g)=0$, hence $g=1$, since
every integer greater than one has a positive term in~\eqref{eq:rate}.
That would contradict admissibility on inputs at least two.
Finally, if a reduced rational $a/b$ makes $(a/b)F$ integral, $b$
divides every coordinate of $F$, so its gcd forces $b=1$.
\end{proof}

Ufnarovski and \AA hlander's Theorem 18 gives a necessary and
sufficient squarefree-denominator condition for nonzero rational
solutions of the logarithmic equation and describes all such solutions
from a particular one~\cite[Section 8]{ufnarovski}. Integer solutions
require nonnegative prime exponents. Kovi\v{c}'s Proposition 2
classifies the positive integer solutions at integer rates~\cite{kovic}.
The following formulation gives the exact existence criterion and all
positive integer solutions at nonnegative rational rates, including
the additional condition $k\ge0$. The same formulation and the
rate-$1/6$ obstruction already appear in the author's earlier
fixed-form work~\cite[Section 3]{corti5}; they are recalled here for
the depth-eight application, without a new general-theorem claim.

\begin{theorem}[All positive integer scales]\label{thm:core}
Let $\lambda=A/B\ge0$ be reduced, $B>0$.
If $B$ is not squarefree, $\Lrate(s)=\lambda$ has no positive integer
solution. If $B$ is squarefree, define, for every prime $p\mid B$,
\[
 c_p\in\{1,\ldots,p-1\},\qquad
 c_p\equiv A(B/p)^{-1}\pmod p,
\]
and put
\begin{equation}\label{eq:core}
 s_0=\prod_{p\mid B}p^{c_p},\qquad
 \alpha_0=\sum_{p\mid B}\frac{c_p}{p},\qquad k=\lambda-\alpha_0.
\end{equation}
For $B=1$, take $s_0=1$ and $\alpha_0=0$. Then $k$ is an integer.
A positive integer solution exists exactly when $k\ge0$, and all
solutions are
\begin{equation}\label{eq:allscales}
 s=s_0\prod_{\ell=1}^{k}p_\ell^{p_\ell},\qquad p_\ell\in\PP.
\end{equation}
Primes may repeat and may divide $B$.
\end{theorem}
\begin{proof}
For any $s=\prod p^{e_p}$, the denominator of $\Lrate(s)=\sum e_p/p$
divides a squarefree product; this proves the first assertion.
Assume $B$ squarefree. The integer
$Bk=A-\sum_{p\mid B}c_p(B/p)$ is divisible by every prime of $B$,
by the defining congruences. It is thus divisible by $B$, so $k$ is
integral, whether or not a solution already exists.

For a solution, let $M$ be the product of the distinct primes dividing
$sB$. Multiplication of the rate equation by $M$ gives
\[
 \sum_{p\mid s}e_p M/p=A M/B.
\]
A prime $p\mid B$ must divide $s$: otherwise reduction modulo $p$
would make the left side zero and the right side nonzero.
Reducing modulo such a $p$ and cancelling the unit $M/B$ gives
$e_p\equiv c_p\pmod p$. For a prime $p\mid s$ outside $B$,
reduction instead gives $e_p\equiv0\pmod p$.
Therefore
$e_p=c_p+ph_p$ on $B$, and $e_p=ph_p$ outside $B$, with
all $h_p\ge0$. Substitution gives $k=\sum h_p\ge0$.
Listing a prime $p$ exactly $h_p$ times produces
\eqref{eq:allscales}. Conversely, every product in that equation has
rate $\alpha_0+k=\lambda$, proving sufficiency and completeness.
\end{proof}

For example, $\lambda=1/6$ has residue core $s_0=2\cdot3^2=18$,
$\alpha_0=7/6$, and $k=-1$: its squarefree denominator does not
give an integer solution. Rate $1/2$ has only scale two, whereas rate
$3/2$ has exactly the scales $2p^p$. In general $k=0$ gives a singleton
scale set; $k\ge1$ gives infinitely many scales by fixing all but one
prime parameter. The rates used below are only zero and one, whose
scale sets are respectively $\{1\}$ and $\{p^p:p\in\PP\}$.
The latter recovers the classical positive fixed-point classification.

\section{The five observed forms}\label{sec:forms}

Define the two primitive forms
\begin{equation}\label{eq:infiniteforms}
\begin{split}
 F_0={}&(1,1,1,1,1,1,1,1,1,3),\\
 F_1={}&(330004229,330099083,341481839,361898942,\\
       &\quad543260957,560785535,672942647,745727049,\\
       &\quad1286744865,1933839382).
\end{split}
\end{equation}
The subscripts of $F_0$ and $F_1$ are labels; both forms have rate
one. Only the other three forms carry observed sequence indices as
subscripts. They are denoted $E_i$, distinct from the depth sets $S_h$:
\begin{equation}\label{eq:singleforms}
\begin{split}
 E_4={}&(398747,1692,2856,5812,5816,8732,9116,9500,\\
       &\quad15700,399467),\\
 E_6={}&(1072281223,85244735,27198937,1431542,920291,\\
       &\quad1,0,0,0,1076009463),\\
 E_9={}&(10361580899,280042764,373531688,560459372,\\
       &\quad611925640,1040273628,1387031516,1515544356,\\
       &\quad2118671140,10865566899).
\end{split}
\end{equation}
All have gcd one when all ten coordinates are included.
Table~\ref{tab:terms} reconstructs every witness as $gF$ from these
five displayed forms. Thus all positive data needed for the mathematical
analysis are contained in the paper.

\begin{table}[htbp]
\centering\small
\begin{tabular}{rrrrc}
\toprule
$i$ & $a(i)$ & $R$ & $g$ & Form\\
\midrule
1&4&12&4&$F_0$\\
2&27&81&27&$F_0$\\
3&3125&9375&3125&$F_0$\\
4&398747&399467&1&$E_4$\\
5&823543&2470629&823543&$F_0$\\
6&1072281223&1076009463&1&$E_6$\\
7&1320016916&7735357528&4&$F_1$\\
8&8910114183&52213663314&27&$F_1$\\
9&10361580899&10865566899&1&$E_9$\\
\bottomrule
\end{tabular}
\caption{The nine-term local prefix and positive roots. The full ordered
witness is the indicated primitive form multiplied by its ten-coordinate
gcd $g$. The first five ranks are inherited; ranks six through nine have
the coverage dependencies in Theorem~\ref{thm:result}.}
\label{tab:terms}
\end{table}

\begin{proposition}\label{prop:classes}
The complete homothetic classes of $F_0$ and $F_1$ have first
coordinates $p^p$ and $330004229p^p$, respectively, for $p\in\PP$.
Each $E_i$, $i=4,6,9$, has exactly one witness in its homothetic
class, namely itself. The nine witnesses therefore represent exactly
five forms: two infinite classes and three singleton classes.
\end{proposition}
\begin{proof}
For $F_0$, $\D(1)=0$ gives rate one, all transitions are
$T_1(1)=1$, and the nine squares sum to $9=3^2$.
For $F_1$, let $u_j$ denote its chain coordinates and let
$c=u_0=330004229=3617\cdot91237$.
The exact derivatives in Table~\ref{tab:seed} satisfy
$u_{j+1}=u_j+\D(u_j)$ for every $0\le j<8$, and
\begin{equation}\label{eq:seedsquare}
 \sum_{j=0}^8u_j^2=3739734755374141924=1933839382^2.
\end{equation}
The gcd of the ten coordinates is one. These identities and
Lemma~\ref{lem:scale} show that $p^pF_1$ is a witness for every
prime $p$, with no exception when $p$ divides $c$ or another
coordinate. Conversely, its first transition forces $\Lrate(s)=1$,
and Theorem~\ref{thm:core} gives exactly $s=p^p$.

Each $E_i$ is an ordinary derivative chain, and its nine chain
coordinates have square sum equal to its last coordinate squared;
these finite identities follow by direct reconstruction from
\eqref{eq:derivative} and~\eqref{eq:singleforms}. Its first transition gives $\lambda=0$.
Equation~\eqref{eq:rate} then forces $s=1$, since every integer
$s>1$ has strictly positive logarithmic derivative.
The five forms displayed are pairwise distinct; normalizing the nine
rows gives precisely the associations of Table~\ref{tab:terms}.
\end{proof}

\begin{table}[htbp]
\centering\small
\begin{tabular}{rrrr}
\toprule
$j$ & $u_j$ & $\D(u_j)$ & $u_{j+1}$\\
\midrule
0&330004229&94854&330099083\\
1&330099083&11382756&341481839\\
2&341481839&20417103&361898942\\
3&361898942&181362015&543260957\\
4&543260957&17524578&560785535\\
5&560785535&112157112&672942647\\
6&672942647&72784402&745727049\\
7&745727049&541017816&1286744865\\
\bottomrule
\end{tabular}
\caption{All eight exact transitions for $F_1$: the last column is the
sum of the middle two columns. These are generalized $T_1$ transitions,
not ordinary derivative transitions of the primitive tuple.}
\label{tab:seed}
\end{table}

The two observed specializations are
\[
 330004229\cdot2^2=1320016916,
 \qquad330004229\cdot3^3=8910114183.
\]
The next prime specialization is
$330004229\cdot5^5=1031263215625$. It is an exact unindexed
witness outside the scanned interval, and is larger than the
historical stationary witness $11^{11}$. Hence it does not improve
the upper bound for the next rank. Since $p^p$ increases with $p$,
the two observed values exhaust this particular class through the
reported frontier. Likewise, $p=2,3,5,7$ exhaust $F_0$ there.

The singleton assertion is global for each specified ordered form:
no other positive integer scale produces that same form. It does not
assert that these integers cannot lie in a parameterized family whose
form changes with the parameter. Nor does a complete list of forms in
a finite prefix prove that no additional forms exist.

\section{Zero tails, depth relations and dynamics}\label{sec:depth}

For a fixed point $p^p$, the sum at depth $h$ is
\begin{equation}\label{eq:stationarydepth}
 Q_h(p^p)=(h+1)p^{2p}.
\end{equation}
It is a square exactly when $h+1$ is a square. This explains the
roles of four legs at depth three and nine legs at depth eight;
fixed points are not members at every depth.

The related OEIS square-sum sequences A230543, A249105, A249106,
and A249107 correspond to depths $2,3,4,5$, respectively
\cite{depthentries}. A210503 adds primitivity at depth one
\cite{triples}; the twin-prime construction of primitive triples is
also discussed by Balzarotti and Lava~\cite[Chapter 22]{book}.
That restriction is absent from $S_8$. Individual higher-derivative
sequences encode chain coordinates rather than ranks: A258651 is
the iterated-derivative array, and A258648 its depth-eight column
\cite{arrayentries}.

\begin{proposition}[Persistence of the known terminating witness]
\label{prop:tail}
Let $N=1072281223$ and $R=1076009463$. Then
\[
 \D^5(N)=1,\qquad\D^j(N)=0\ (j\ge6),\qquad
 Q_h(N)=R^2\ (h\ge5).
\]
Moreover $N\notin S_4$.
\end{proposition}
\begin{proof}
The exact chain is the nine chain coordinates of $E_6$ in
\eqref{eq:singleforms}. Here $x_5=\D^5(N)=1$, and
\eqref{eq:derivative} supplies the zero tail. Its square sum through
depth five is $R^2=1157796364465548369$; adding zero squares changes
nothing. At depth four the sum is $R^2-1$, which lies strictly
between $(R-1)^2$ and $R^2$.
\end{proof}
This persistence, including its depth-eight consequence, was already
established in the A249107 study~\cite{corti7}. There $N$ has rank
twelve in the author's extended A249107 sequence; the public OEIS
entry checked on October 27, 2026 still has ten terms and does not
include $N$~\cite{depthentries}. Its A249110 rank is six; neither
sequence's rank follows from the other. More generally, if $\D^k(n)=0$, $Q_h(n)$ is constant
for $h\ge k-1$, but it need not be square. A zero tail alone is not
a membership test.

There is no general inclusion of the sets as depth increases. For
example $Q_3(4)=64$ is square and $Q_4(4)=80$ is not. In the other
direction, the chain beginning
$(19164,25564,31848,58412,61916)$ gives
$Q_4(19164)=96336^2$, while
$73804^2<Q_3(19164)=5447033840<73805^2$.
Among the nine inputs of Table~\ref{tab:terms}, the square depths
from one through eight are exactly $\{3,8\}$ for the four
stationary inputs, $\{5,6,7,8\}$ for $N$, and $\{8\}$ for the
remaining four inputs. This is a finite correlation check on those
nine inputs, not an enumeration at any other depth. In particular,
the exclusions of depths one, two and three for $N$ are finite
integer-square-root checks; Proposition~\ref{prop:tail} establishes
the depth-four exclusion and persistence at all depths at least five.

\begin{proposition}[Dynamical types in the observed prefix]\label{prop:dynamics}
The four members of form $F_0$ are stationary, $N$ eventually reaches
zero, and the four other displayed members have derivatives tending
to infinity.
\end{proposition}
\begin{proof}
The first two assertions follow from fixed points and
Proposition~\ref{prop:tail}. For $F_1$, the entire integral
$T_1$ orbit of $c>1$ satisfies $T_1(m)=m+\D(m)\ge m+1$
for $m\ge2$. Lemma~\ref{lem:scale} at every depth therefore shows
divergence for all its witnesses.
For $E_4$ and $E_9$, the first derivatives are respectively
$1692$ and $280042764$, both divisible by four and greater than four.
The classical four-divisor argument~\cite[Theorem 4]{ufnarovski}
can be seen directly: for $m\ge2$,
\[
 \D(4m)=4m+4\D(m)\ge4m+4.
\]
The next value remains a multiple of four, so induction gives
unbounded growth.
\end{proof}
Thus singleton scale classes need not be terminating: two diverge and
one terminates in this sample. These pointwise proofs do not resolve
the general dynamics of the arithmetic derivative.

\section{Exact enumeration and numerical domain}\label{sec:method}

\subsection{First-derivative sieve and the remaining chain}

The producer is a specialized C17 program with GCC's unsigned
128-bit extension~\cite{gcc}, OpenMP, GMP and OpenSSL libcrypto.
It supports original inputs through $50000000000$, a technical
limit distinct from the executed endpoint. The source and build used
for the reported runs underwent external AI-assisted code review
before those runs, as described in the AI use disclosure. This was
neither independent human code review nor a second enumeration of
the negatives; manuscript review is a separate stage.

In a block $[L,H]$ of at most $262144$ inputs, initialize a residual
to each original input $n$ and a derivative accumulator to zero.
Generate all primes through $\lfloor\sqrt H\rfloor$. For each prime
$p$, visit its multiples, remove its complete exponent $e$ from
the residual and add $e(n/p)$, using the original $n$. A final
residual $q>1$ is prime and contributes $n/q$. Indeed after all
primes through $\sqrt H$ have been removed, a composite residual
would have a prime divisor at most $\sqrt H$, a contradiction.
The accumulator is therefore exactly $\D(n)$.

If $\D(n)=1$, $n$ is prime: a composite $n=ab$, $a,b\ge2$, has
$\D(n)=a\D(b)+b\D(a)\ge a+b\ge4$.
A prime input has chain $(n,1,0,\ldots,0)$ and
$n^2<Q_8(n)=n^2+1<(n+1)^2$, so this rejection is exact.
No family or primitivity restriction is used to remove other inputs.

For the remaining inputs, the producer computes the seven further
derivatives by factoring $x_1,\ldots,x_7$ as needed. Small-prime
division precedes a fixed-seed rho accelerator using Pollard's
quadratic map with Floyd cycle detection~\cite{pollard}. Proper factors recurse on strictly smaller
operands. The bounded accelerator uses at most 32 seeds and
$100000$ steps per seed, then complete odd trial division.
Factors are checked as exact divisors, prime leaves are checked,
and their product must reconstruct the operand before the derivative
is evaluated. The fallback is finite, but no useful uniform
practical running-time bound is asserted. Interruption or failure
leaves the block incomplete; it never becomes a negative answer.
Known zeros are filled with zeros, and an established equality of
successive iterates allows the remaining fixed-point coordinates to
be filled exactly.

Every coordinate is imported into GMP before squaring. The nine
squares are added in multiple precision; the program applies
\texttt{mpz\_perfect\_square\_p}, extracts the integer root, and
checks its square again~\cite{gmp}. It does not square a large
coordinate in a narrower native type first.

\subsection{Uniform bounds and the two primality domains}

A bound for all possible intermediates is needed even when observed
positive chains are small. The smallest-prime-factor bound and
two-adic decomposition used here are described in
\cite[Theorems 2.1.1 and 3.2.2]{dahl}, using the classical
derivative-bound argument of~\cite[Theorem 9]{ufnarovski}.
The following integer version
makes the rounding explicit. For $U\ge2$ and
$0\le a\le\lfloor\log_2 U\rfloor$, let $M_a$ be the largest odd
integer at most $\lfloor U/2^a\rfloor$, and put
$k_a=\lfloor\log_3M_a\rfloor$. Define
\begin{equation}\label{eq:majorant}
 G(U)=\max_a\left\{
 a2^{a-1}M_a+2^a\left\lfloor\frac{M_a k_a}{3}\right\rfloor
 \right\},\qquad G(0)=G(1)=0,
\end{equation}
with the first summand zero for $a=0$.

\begin{lemma}\label{lem:bound}
For $0\le n\le U$, $\D(n)\le G(U)$, and $G$ is nondecreasing.
\end{lemma}
\begin{proof}
Write $n=2^a m$, $m$ odd. For $m>1$, each prime factor counted
with multiplicity is at least three, so
$\Omega(m)\le\lfloor\log_3m\rfloor$ and
$\D(m)\le\lfloor m\lfloor\log_3m\rfloor/3\rfloor$.
The inequality also holds for $m=1$.
The product rule gives
$\D(n)=a2^{a-1}m+2^a\D(m)$.
For fixed $a$ its upper expression is nondecreasing in odd $m$;
replace $m$ by $M_a$ and maximize. As $U$ increases each existing
$M_a$ is nondecreasing and only additional choices of $a$ enter,
proving monotonicity. The cases zero and one are separate.
\end{proof}

Starting at $U_0=50000000000$, Lemma~\ref{lem:bound} gives the
induction $x_j\le U_j\Rightarrow x_{j+1}=\D(x_j)\le G(U_j)=U_{j+1}$.
The resulting chain of uniform bounds is Table~\ref{tab:bounds}. The floors can be computed by integer
division, without floating-point logarithms.

\begin{table}[htbp]
\centering\small
\begin{tabular}{rrrc}
\toprule
$j$ & $U_j$ & Bits & Coordinate storage\\
\midrule
0&50000000000&36&unsigned 64-bit\\
1&797790175232&40&unsigned 64-bit\\
2&14345190768640&44&unsigned 64-bit\\
3&294669116243968&49&unsigned 64-bit\\
4&6755399441055744&53&unsigned 64-bit\\
5&174514485560606720&58&unsigned 64-bit\\
6&4643211215818981376&63&unsigned 64-bit\\
7&142962266571249025024&67&unsigned 128-bit\\
8&4574792530279968800768&72&unsigned 128-bit\\
\bottomrule
\end{tabular}
\caption{Universal upper bounds at the supported input cap $5\cdot10^{10}$.
They are not observed maxima and do not assert coverage to that cap.}
\label{tab:bounds}
\end{table}

The sum bound is
$20949186494680010958899369285145516085084160$, requiring 144 bits,
and its integer-root bound is $4577028129111728646739$, requiring
72 bits. Thus the full arithmetic chain is not confined to 64 bits,
and a 128-bit square sum is not justified by these bounds.
Even the single-coordinate bound $U_8^2$ requires 144 bits, so
importing the coordinates into GMP must precede multiplication.

The primality predicate is the strong probable-prime test
of Miller and Rabin~\cite{miller,rabin}, made deterministic on the
admitted finite domains
by two external premises. For operands below $2^{64}$, the bases are
\[
 2,\ 325,\ 9375,\ 28178,\ 450775,\ 9780504,\ 1795265022,
\]
the Jim Sinclair set recorded in~\cite{bases}. Bases are reduced
modulo the operand, with zero residues skipped. The word backend
uses 128-bit modular products. For larger operands, the program
uses the first twelve prime bases, $2,3,\ldots,37$.
Sorenson and Webster establish that the least composite passing
all those tests is
\begin{equation}\label{eq:psi12}
 \psi_{12}=318665857834031151167461
\end{equation}
\cite[Theorem 1.1 of the preprint]{sorenson}.
Only $x_7$ among the factored operands can exceed 64 bits; it is at
most $U_7<\psi_{12}$. The wide test is guarded strictly below
$\psi_{12}$ itself. Its modular multiplication and exponentiation use
GMP; native 128-bit products would not cover every product of
two 67-bit residues. The producer uses the word helper whenever
the operand fits and dispatches to the wide helper only when needed.

These finite-domain premises are cited external computational results,
not new exhaustive primality computations performed in this study.
The negative census retains them as dependencies. Trial division of
the reported positive chains does not use these premises.

\subsection{Parallel ownership, coverage and recovery}

The first-derivative sieve is followed by an OpenMP loop with
dynamic chunks of 64 candidates. Each worker owns its GMP scratch
and writes a separate result flag. The implicit join precedes the
serial examination of flags and replay of positive inputs in
increasing order~\cite{openmp}. A block is committed only if all
its candidates have completed. The program's statuses distinguish
positive, negative and incomplete evaluations. A signal or resource
failure prevents advancement past the incomplete block.

To specify the partition, let $w=262144$ and let
$K=\lceil(H-L+1)/w\rceil$ be the number of blocks in a configured
interval. With $\sigma$ operational segments, segment $i$, $0\le i<\sigma$,
uses block indices
$\lfloor iK/\sigma\rfloor$ through
$\lfloor(i+1)K/\sigma\rfloor-1$.
Clamping the last upper endpoint to $H$ gives a disjoint consecutive
partition. The state records configuration identity, segment bounds,
completed frontier, block count, saved work time, fallback count,
and ordered positive tuples. Every completed segment receives a
seal. Consolidated manifests and positive tables are checked against
those authoritative states; a future unstarted segment records no
examined interval. Integrity records detect changed bytes and bind
the recorded producer; they do not by themselves prove that each
negative classification was mathematically correct.

\section{The finite result and its validation}\label{sec:results}

\begin{theorem}[Computational extension, with explicit dependencies]
\label{thm:result}
Subject to the accepted five-term baseline, correct execution of the
AI code-reviewed enumeration described in Section~\ref{sec:method}, its
stated numerical and finite-domain primality premises, and the
complete saved records, the only members of $S_8$ in the inclusive
interval $823544\le n\le22000000000$ are
\[
 1072281223,\quad1320016916,\quad8910114183,\quad10361580899.
\]
They have consecutive ranks six through nine. No further member
occurs after $a(9)$ through the inclusive endpoint $22000000000$.
\end{theorem}
\begin{proof}[Proof (computational)]
The classified intervals are displayed in Table~\ref{tab:coverage}.
The lower endpoint of each continuation is exactly one greater than
the preceding upper endpoint. There are 34 sealed segments and
$83921$ committed blocks. Their input counts sum to
\[
 1071457680+19174260736+1753458041
   =21999176457=22000000000-823543.
\]
The complete-block rule and checked partitions supply the recorded
coverage, under the producer assumptions in the statement.
The first interval has only $N$, exactly at its upper endpoint.
The second has the other three positives in increasing order and
the third has none. Their full chains agree with the documentary
table. Appending these four hits to the inherited prefix assigns
the displayed ranks. Independent positive reconstruction checks
the eight derivative identities and square identity for every row;
the exclusion component rests on the producer and records, not on
that positive-only check.
\end{proof}

\begin{table}[htbp]
\centering\small
\setlength{\tabcolsep}{4pt}
\begin{tabular}{lrrrr}
\toprule
Portion used & Inclusive lower & Inclusive upper & Seals & Hits\\
\midrule
First run&823544&1072281223&2&1\\
Continuation prefix&1072281224&20246541959&29&3\\
Final run&20246541960&22000000000&3&0\\
\bottomrule
\end{tabular}
\caption{Executed coverage after the inherited five-term prefix.
The continuation was configured for $50000000000$ and remains
incomplete relative to that original target. Only its sealed
29-segment prefix contributes. The final run is separately complete.}
\label{tab:coverage}
\end{table}

\begin{corollary}\label{cor:next}
With the computational dependencies of Theorem~\ref{thm:result},
\begin{equation}\label{eq:next}
 22000000000<a(10)\le285311670611=11^{11}.
\end{equation}
The value of $a(10)$ is unknown.
\end{corollary}
\begin{proof}
The theorem gives the strict lower bound. The known stationary
family supplies $11^{11}\in S_8$ above the frontier. Infinitely
many prime parameters supply infinitely many distinct members,
so the next rank exists. The witness $11^{11}$ need not have rank
ten; it supplies only the upper bound.
\end{proof}

The original continuation configuration was preserved when the
author selected the narrower endpoint. The final adjacent run has
its own configuration; no old target or unstarted state was
relabelled as completed coverage. The supported input cap remains
$5\cdot10^{10}$, which also lies below the mathematical upper
witness. Capacity is not evidence that those additional inputs
were examined.

\subsection{What the independent checks establish}\label{sec:checks}

The independent Python record checker imports no producer code and
invokes no producer. It checks configuration and state grammar,
frontiers, seals, ordered joins, counts and correspondence with the
consolidated records. For every positive tuple, it independently
factors by complete trial division, reconstructs the chain, and
uses integer square roots. The documentary-table checker compares
the indexed and witness tables row by row, checking all 72
derivative transitions and all nine square identities.
Python arbitrary-size integers and \texttt{math.isqrt}
\cite{python} give an arithmetic path distinct from the C
factorization and GMP square test.

These checks confirm saved positive arithmetic and record
consistency. They do not independently visit every negative input
in the added interval, and they share the mathematical definition
and saved inputs with the producer. Multiple Python implementations
of this arithmetic are corroboration of the same positive axis,
not multiple exhaustive scans. The original first five ranks and
their gaps remain inherited, even though their membership is
rechecked. The retained finite checker for this paper additionally
normalizes the nine tuples, checks every generalized transition,
reconstructs the majorant constants and manifest totals, and tests
the stated depth correlations on these nine inputs.

The unchanged development evidence includes small direct comparison
windows, known positives, boundary/refusal cases, forced fallback,
signal/recovery and parser tests, and word/wide arithmetic fixtures.
Some historical development drivers require the original development
tree and dependencies and are not replay commands for the review
extract; the independent record, documentary-table and finite
structural checkers described above are replayable.
The wide fixtures include recursive Lucas $n-1$ prime certificates
with checked factor products; these are scoped checks of selected
operands. A maximum-width transport fixture was deliberately not
a mathematical witness and was rejected by the mathematical
verifiers. It contributes no sequence term. Recorded GCC/Clang,
sanitizer, static-analysis and qualified OpenMP race-detector checks
support the implementation obligations in their tested domains;
none is a mathematical certificate for all negatives.

\subsection{Recorded environment and measurement scope}

The frozen environment records an AMD Ryzen 9 7940HS with eight
physical cores and sixteen logical processors, Linux Mint 22.3,
kernel 6.17.0-42, GCC 13.3.0, GMP 6.3.0 and OpenSSL 3.0.13.
The runs use sixteen workers and record OpenMP macro value 201511.
The build uses \texttt{-O3 -std=c17 -fopenmp}, strict warning
options promoted to errors, and links GMP and libcrypto, without
processor-specific tuning flags.

The sum of saved completed-block work counters is approximately
$26595$ seconds, about 7 hours 23 minutes: approximately $914$,
$23332$, and $2349$ seconds for the three portions respectively.
These counters time completed block work, not complete campaign
wall time. They exclude checkpoint and audit overhead, pauses,
and any repeated unfinished work. The saved fallback-input count
is zero; that observation does not remove the fallback contract
or bound the cost of future inputs. No sustained campaign peak
RSS, installed-memory capacity or thermal qualification is inferred
from short development measurements.

\subsection{Local reproducibility and publication boundary}

The local project retains the source and build, original executable,
canonical indexed and documentary tables, configurations, sealed
states, manifests and runnable Python validators. Auditing the
saved records and reconstructing the positive witnesses requires
no new official enumeration. A complete public reproduction also
requires those actual run payloads and their configuration bindings;
the manuscript's tables alone cannot reproduce the negative census.
The project reference~\cite{project} identifies the public repository
and Zenodo release series. The initial source-and-evidence package is
available as a GitHub release and as an open Zenodo deposit.
The accompanying package provides runnable retained-data, positive-witness
and structural checks. Those checks audit the saved evidence; they do
not independently repeat the full negative census.

\section{Consequences and open questions}\label{sec:questions}

The finite prefix combines four stationary examples, two members
of one nonconstant infinite form, and three singleton forms.
Thus two different distinctions are necessary: membership versus
rank, and complete scales of a fixed form versus a classification
of all forms. A new family member gives an upper witness, while
its rank requires intervening exclusions. Likewise, the rate
criterion completely determines the scales of $F_1$ but supplies
no complete enumeration of other primitive shapes. Any further
member strictly below $11^{11}$ must have a primitive form absent
from Table~\ref{tab:terms}: the next $F_0$ member is $11^{11}$,
the next $F_1$ member is larger, and the three $E_i$ classes have
no further scales. This does not force $a(10)$ to have a new form,
since equality with the upper witness remains possible.

The general rate theory permits rational rates, subject to
Theorem~\ref{thm:core}. All five observed forms have rates zero
or one. Whether an admissible depth-eight form has a noninteger
rate, whether infinitely many primitive forms occur, and whether
infinitely many of them are singleton or terminating forms remain
unresolved here. So does the least member above the completed
frontier. The nine observations provide no justified density law,
probability model for a future hit, or asymptotic estimate.
The two infinite classes and the three specified singleton classes,
in contrast, follow by proof rather than extrapolation.

\section{Conclusion}\label{sec:conclusion}

The completed added coverage extends the accepted five-term
prefix to nine consecutive terms and excludes further members
through $22000000000$, with explicit production and baseline
dependencies. Its first added rank belongs to a known persistent
zero-tail witness. The ordered forms separate two infinite
homothetic classes from three singleton classes, and the
nonconstant class explains the proportional tuples at ranks seven
and eight. These structural results account for the observed
prefix without classifying all possible forms. The next rank is
bounded by a historical fixed-point witness and requires further
minimality evidence for its determination.

\section*{Data availability}

The source and build, original configurations and sealed records,
indexed b-file, documentary derivative/root table, independent
positive and record checkers, finite structural checks, environment
captures and validation summaries are available from the public
GitHub repository and Zenodo release series cited in the References.
The package also includes this monolithic LaTeX source and its
matching PDF. The repository documents the verification commands,
dependencies and restoration of the original evidence layout.
The Zenodo DOI identifies the release series, rather than a single
immutable version; individual deposits retain their own version-specific
DOIs. Positive reconstruction and record audits do not independently
certify all negatives in the reported interval.

\section*{AI use disclosure}

OpenAI Codex assisted with research, implementation work, local
verification and preparation of this manuscript. The supplied code
review reports attest reported contributions from Anthropic Claude,
GLM/Z.ai, xAI Grok, Meta AI, MiniMax and Qwen, alongside OpenAI.
These are reported families rather than independently authenticated
backend identities; exact versions are not consistently established.
A subsequent manuscript round supplied ten AI-attributed reports.
Their comments were adjudicated against the manuscript and primary
sources; model identities and reported executions were assessed
separately from the evidence actually supplied.

\paragraph{Mathematical analysis.}
Codex assisted in reconciling the OEIS definition and historical
witnesses, developing the product-rule and fixed-form arguments,
checking the nine normalized tuples and investigating depth
relations. The mathematical statements are supported by the
displayed proofs, primary source readings and exact finite
calculations, with historical principles attributed to their sources.

\paragraph{Algorithm and code.}
Codex assisted with the C implementation and its numerical,
parallel and recovery contracts. The reported external families
above contributed code review; their reports were adjudicated
against the source and retained tests. The author executed the official
campaigns. Codex executed the retained Python record and positive
checks locally during manuscript preparation and revision. The numerical result is tied to that producer's records,
known-term checks, separate positive arithmetic and scoped
development tests. Code-review agreement is not an independent
negative enumeration or manuscript review.

\paragraph{Manuscript.}
Codex assisted in English and LaTeX drafting, literature checking
and internal review of proofs, numerical statements, chronology,
citations and availability. Final numerical claims rely on local
artifacts and validation, not conversational assertions. The new
manuscript proofs and the adjudicated review findings remain
subject to the author's scientific assessment. The recorded checker
executions identify their operator; they do not attest to a personal
rerun of every check by the author. Carlo Corti retains full scientific
responsibility for the mathematical statements, algorithm,
implementation, numerical results, validation boundaries and
conclusions.

\clearpage
\begin{thebibliography}{99}
\sloppy\raggedright

\bibitem{sequence}
P. P. Lava, OEIS A249110, Numbers that form a Pythagorean 10-tuple
with their first eight arithmetic derivatives, October 21, 2014;
subsequence contribution by R. Chandler, December 22, 2016.
\url{https://oeis.org/A249110}.
Revision 24 snapshot acquired October 24, 2026; entry checked
October 27, 2026.

\bibitem{barbeau}
E. J. Barbeau, Remarks on an Arithmetic Derivative,
\emph{Canadian Mathematical Bulletin} \textbf{4} (1961), no. 2, 117--122.
DOI: \texttt{10.4153/CMB-1961-013-0}.
\url{https://www.cambridge.org/core/services/aop-cambridge-core/content/view/1FD7F09AD3972692FC97BB23A21D0BD8/S0008439500050773a.pdf/remarks_on_an_arithmetic_derivative.pdf}.

\bibitem{ufnarovski}
V. Ufnarovski and B. \AA hlander, How to Differentiate a Number,
\emph{Journal of Integer Sequences} \textbf{6} (2003), Article 03.3.4.
\url{https://cs.uwaterloo.ca/journals/JIS/VOL6/Ufnarovski/ufnarovski.html}.

\bibitem{kovic}
J. Kovi\v{c}, The Arithmetic Derivative and Antiderivative,
\emph{Journal of Integer Sequences} \textbf{15} (2012), Article 12.3.8.
\url{https://cs.uwaterloo.ca/journals/JIS/VOL15/Kovic/kovic4.html}.

\bibitem{fixedpoints}
A. Auel, OEIS A051674, $a(n)=\operatorname{prime}(n)^{\operatorname{prime}(n)}$.
Fixed-point comment by T. D. Noe, October 12, 2004.
\url{https://oeis.org/A051674}. Consulted October 27, 2026.

\bibitem{corti2}
C. Corti, \emph{An Exact Local Extension and Homothetic Classes for
OEIS A230543}, 2026, Sections 2--3.
Release-series DOI: \texttt{10.5281/zenodo.23311392}.
\url{https://github.com/carcorti/A230543/blob/v1.0.1/paper/A230543.tex}.
Release v1.0.1 manuscript consulted October 27, 2026.

\bibitem{corti6}
C. Corti, \emph{An Exact Extension and an Infinite Family for
OEIS A249106}, 2026, Sections 3--4.
Release-series DOI: \texttt{10.5281/zenodo.23155533}.
\url{https://github.com/carcorti/A249106/blob/v1.0.1/paper/A249106.tex}.
Release v1.0.1 manuscript consulted October 27, 2026.

\bibitem{corti5}
C. Corti, \emph{A Consecutive Extension and Fixed-Form Families for
OEIS A249105}, 2026, Sections 2--3.
Release-series DOI: \texttt{10.5281/zenodo.23388128}.
\url{https://github.com/carcorti/A249105/blob/v1.0.1/paper/A249105.tex}.
Release v1.0.1 manuscript consulted October 27, 2026.

\bibitem{corti7}
C. Corti, \emph{Arithmetic-Derivative Pythagorean 7-Tuples: Exact
Extension and Homothetic Structure of OEIS A249107}, 2026,
Sections 2 and 4.
Release-series DOI: \texttt{10.5281/zenodo.23534992}.
\url{https://github.com/carcorti/A249107/blob/v1.0.1/paper/A249107.tex}.
Release v1.0.1 manuscript consulted October 27, 2026.

\bibitem{depthentries}
P. P. Lava, OEIS A230543 (October 25, 2013), A249105, A249106 and
A249107 (October 21, 2014), arithmetic-derivative square-sum
sequences at depths 2, 3, 4 and 5. Extensions include G. Resta,
R. Chandler and C. Corti as credited in the respective entries.
\url{https://oeis.org/A230543}; \url{https://oeis.org/A249105};
\url{https://oeis.org/A249106}; \url{https://oeis.org/A249107}.
Project snapshots consulted October 27, 2026.

\bibitem{triples}
P. P. Lava, OEIS A210503, Numbers forming a primitive Pythagorean
triple with their arithmetic derivative, January 25, 2013;
further comments and table by R. Chandler.
\url{https://oeis.org/A210503}. Consulted October 27, 2026.

\bibitem{book}
G. Balzarotti and P. P. Lava, \emph{La derivata aritmetica.
Alla scoperta di un nuovo approccio alla teoria dei numeri},
Hoepli, Milan, 2013. Print ISBN 978-88-203-5864-8;
digital edition ISBN 978-88-203-6038-2.
Targeted chapters and appendices consulted; not an integral reading.

\bibitem{arrayentries}
A. P. Heinz, OEIS A258651, Array of iterated arithmetic derivatives,
and A258648, Eighth arithmetic derivative of $n$, June 6, 2015.
\url{https://oeis.org/A258651}; \url{https://oeis.org/A258648}.
Project snapshots consulted October 27, 2026.

\bibitem{dahl}
N. Dahl, J. Olsson and A. Loiko, \emph{Investigations on the
properties of the arithmetic derivative}, arXiv:1108.4762v1, 2011.
Preprint DOI: \url{https://doi.org/10.48550/arXiv.1108.4762}.

\bibitem{pollard}
J. M. Pollard, A Monte Carlo method for factorization,
\emph{BIT} \textbf{15} (1975), no. 3, 331--334.
\url{https://doi.org/10.1007/BF01933667}.

\bibitem{miller}
G. L. Miller, Riemann's hypothesis and tests for primality,
\emph{Journal of Computer and System Sciences} \textbf{13} (1976),
no. 3, 300--317. DOI: \texttt{10.1016/S0022-0000(76)80043-8}.
\url{https://www.cs.cmu.edu/~glmiller/Publications/Papers/Mi76.pdf}.

\bibitem{rabin}
M. O. Rabin, Probabilistic algorithm for testing primality,
\emph{Journal of Number Theory} \textbf{12} (1980), no. 1, 128--138.
DOI: \texttt{10.1016/0022-314X(80)90084-0}.
\url{https://www.sciencedirect.com/science/article/pii/0022314X80900840}.

\bibitem{bases}
W. Izykowski, \emph{Deterministic variants of the Miller--Rabin
primality test: SPRP bases records}, Jim Sinclair's seven-base
record dated April 20, 2011.
\url{https://miller-rabin.appspot.com/}. Consulted October 27, 2026.

\bibitem{sorenson}
J. P. Sorenson and J. Webster, Strong Pseudoprimes to Twelve Prime
Bases, \emph{Mathematics of Computation} \textbf{86} (2017),
no. 304, 985--1003. Journal DOI: \texttt{10.1090/mcom/3134}.
Preprint arXiv:1509.00864v1 (2015),
Theorem 1.1, consulted here.
Preprint DOI: \url{https://doi.org/10.48550/arXiv.1509.00864}.

\bibitem{gmp}
The GMP Development Team, \emph{GNU MP 6.3.0 Manual},
Integer Roots and Integer Import and Export.
\url{https://gmplib.org/manual/Integer-Roots};
\url{https://gmplib.org/manual/Integer-Import-and-Export}.
Consulted October 27, 2026.

\bibitem{gcc}
Free Software Foundation, \emph{Using the GNU Compiler Collection,
GCC 13.3.0}, 128-bit Integers.
\url{https://gcc.gnu.org/onlinedocs/gcc-13.3.0/gcc/_005f_005fint128.html}.
Consulted October 27, 2026.

\bibitem{openmp}
OpenMP Architecture Review Board, \emph{OpenMP Application
Programming Interface}, version 4.5, November 2015, Sections 2.5
and 2.7.1. \url{https://www.openmp.org/wp-content/uploads/openmp-4.5.pdf}.
Consulted October 27, 2026.

\bibitem{python}
Python Software Foundation, \emph{Python 3.12 Documentation},
\texttt{math.isqrt}.
\url{https://docs.python.org/3.12/library/math.html\#math.isqrt}.
Consulted October 27, 2026.

\bibitem{maple}
Maplesoft, \emph{Maple Online Help}, \texttt{ifactors},
examples (4)--(5).
\url{https://www.maplesoft.com/support/help/Maple/view.aspx?path=ifactors}.
Consulted October 27, 2026.

\bibitem{project}
C. Corti, \emph{Nine-Square Arithmetic-Derivative Tuples: An
Extension of OEIS A249110 and Their Scaling Structure},
GitHub repository and Zenodo archive, 2026.
\url{https://github.com/carcorti/A249110};
\url{https://doi.org/10.5281/zenodo.23695423}.
Public repository and release-series archive. The Zenodo DOI is the
series (concept) identifier; version-specific deposits have distinct DOIs.

\end{thebibliography}
\end{document}
